\documentclass{leveldoc}
\DocUses{foundations,first_bend,twin_beams}
\DocFingerprint{34c46d382b3df452}

\begin{document}

\section{What the level asks}\label{sec:asks}

The game's text:
\leveldescription

One particle ($q = \val{q}$, $m = \val{m}$) is launched from $(\val{xLaunch}, \val{yLaunch})$
at $\val{launchAng}^\circ$ above $+x$, with kinetic energy $T_0 = \val{Tzero}$. The detector
is the rectangle $\val{xD} \le x \le \val{xDhi}$, $\val{yDlo} \le y \le \val{yDhi}$. An entry
counts only when the velocity there lies within $\val{halfDeg}^\circ$ of the $+x$ axis. The
player places up to three charges, of either sign and of the magnitudes in the level file,
on any node. A flight ends when the particle touches a charge, leaves the world, enters the
detector or runs out of time (\xsecref{foundations}{sec:board}); $t_{\max} = \val{tmax}$.

The reference solution is the charge $Q = \val{Q}$ at $(\val{xC}, \val{yC})$. The coupling
$\kappa = qQ = \val{kappa}$ is attractive, and the charge sits below the ray, so the ray is
bent down. It meets the detector on the floor, travelling at $\val{angE}^\circ$ to $+x$,
inside the cone. This document cites the orbit of one charge and derives the cone.

\section{Setup}\label{sec:setup}

The level file, field by field (the appendix reproduces it):
\setuptable

In symbols: $c = \val{c}$, the particle has radius $a = 0$, and the charge has radius
$R = \val{R}$. The launch has
\[
  \gamma_0 = 1 + \frac{T_0}{mc^2} = \val{gammaZero}, \qquad
  \abs{\vect{p}_0} = \sqrt{\frac{T_0^2}{c^2} + 2mT_0} = \val{pZero}, \qquad
  v_0 = \frac{\abs{\vect{p}_0}}{\gamma_0 m} = \val{vZero}.
\]
The direction stored in the file makes the angle $\val{launchAng}^\circ$ with $+x$. That
angle and $20^\circ$ differ by $\val{dirResidual}$ radians, and the stored components
differ in length from $1$ by $\val{dirNormGap}$, so the aim used below is the file's. The world is
$-1 \le x \le \val{xWorld}$, $-1 \le y \le \val{yWorld}$. The cone's axis is the angle
$\val{axisDeg}^\circ$ and its half-angle is $\val{halfDeg}^\circ$, so the axis is the unit
vector $\uvec{e} = (1, 0)$ and the half-angle in radians is $\delta = \val{halfRad}$. The
detector sets no condition on the kinetic energy.

\begin{figure}[htbp]
\centering
\begin{tikzpicture}
\begin{axis}[axis equal image, width=0.92\linewidth, xmin=-1.6, xmax=33, ymin=-1.6, ymax=21.6,
  xtick={0,5,...,30}, ytick={0,5,...,20}, tick label style={font=\small},
  xlabel={$x$}, ylabel={$y$}, label style={font=\small}]
  \draw[black!8] (0,0) grid (30,20);
  \draw[black!45, dashed] (-1,-1) rectangle (31,21);
  \fill[green!50!black, fill opacity=0.18] (28,12) rectangle (30,17);
  \draw[green!40!black] (28,12) rectangle (30,17);
  \fill[teal, fill opacity=0.28] (\rawval{coneOx},\rawval{coneOy}) --
    (\rawval{coneXp},\rawval{coneYp}) -- (\rawval{coneXm},\rawval{coneYm}) -- cycle;
  \draw[teal, thick] (\rawval{coneOx},\rawval{coneOy}) -- (\rawval{coneXp},\rawval{coneYp});
  \draw[teal, thick] (\rawval{coneOx},\rawval{coneOy}) -- (\rawval{coneXm},\rawval{coneYm});
  \addplot[thin, black!45] table {pathfree.dat};
  \addplot[thick, orange!80!black] table {pathhalf.dat};
  \addplot[very thick, blue!70!black] table {path.dat};
  \fill[red!70!black] (21,0) circle[radius=0.28];
  \node[font=\small, red!60!black, anchor=south] at (21,0.35) {$Q$};
  \draw[-{Stealth[length=6pt]}, thick] (\rawval{xLaunch},\rawval{yLaunch}) --
    (\rawval{arrowX},\rawval{arrowY});
  \fill (\rawval{xLaunch},\rawval{yLaunch}) circle[radius=0.12];
  \node[font=\small, green!30!black, anchor=south] at (29,17.25) {detector};
  \node[font=\small, teal, anchor=west] at (32.05,14.5) {$\pm 8^\circ$};
\end{axis}
\end{tikzpicture}
\caption{The reference orbit (blue) meets the floor of the detector. The teal wedge is the
cone of accepted directions, $\pm\val{halfDeg}^\circ$ about $+x$, drawn from the middle of
the entrance. The weaker charge (orange; §\ref{sec:others}) and the straight ray (grey)
both enter through the left face, at $\val{halfAng}^\circ$ and $\val{freeAng}^\circ$.}
\label{fig:board}
\end{figure}

\section{Principles and model}\label{sec:model}

\begin{principles}
  \item The field of a charged sphere and the force on a non-polarizable particle,
    \xeqref{foundations}{eq:coulomb} and \xeqref{foundations}{eq:sphereforce}.
  \item The relativistic equation of motion \xeqref{foundations}{eq:eom}, the launch
    \xeqref{foundations}{eq:launch}, the energy \xeqref{foundations}{eq:energy} and the
    angular momentum \xeqref{foundations}{eq:angmom}. With one charge, the orbit
    \xeqref{first_bend}{eq:general}.
  \item The board's rules (\xsecref{foundations}{sec:board}). The first entry into the
    detector decides, and the state at that entry has to lie in the cone of §\ref{sec:cone}.
  \item Newton's method and Simpson's rule (\xsecref{foundations}{sec:numerics}).
\end{principles}

With the charge at $\vect{r}_\mathrm{c}$ the force is
$\kappa(\vect{x}-\vect{r}_\mathrm{c})/\abs{\vect{x}-\vect{r}_\mathrm{c}}^3$, and the state
obeys \xeqref{foundations}{eq:eom} until the first event. The motion stays in the plane
(\xsecref{foundations}{sec:plane}).

\begin{assumption}{The charge is held in place}\label{as:fixed}
The charge does not recoil. That is a rule of the game (\xsecref{foundations}{sec:radiation}).
\end{assumption}

\begin{assumption}{No radiation}\label{as:radiation}
The radiated energy \xeqref{foundations}{eq:larmor} is not taken from the particle. Along
the reference flight it is the fraction $\val{EradFrac}$ of $T_0$, and along the weaker
charge of §\ref{sec:others} the fraction $\val{halfEradFrac}$ (§\ref{sec:hand}). With no
charge the acceleration is zero, so the radiated energy is zero. The game makes the same
approximation and reports the same integral (\physref{8}).
\end{assumption}

\begin{assumption}{A point particle}\label{as:point}
$a = 0$, and the particle does not polarize, so \xeqref{foundations}{eq:sphereforce} is
exact. On the reference flight the closest approach is $\val{rmin}$, which is
$\val{rOverR}$ times $R$. The weaker charge stays at least $\val{halfRmin}$ away.
\end{assumption}

\section{Derivation}\label{sec:derivation}

\subsection{The orbit}\label{sec:orbit}

Take polar coordinates about the charge,
$\vect{x} - \vect{r}_\mathrm{c} = r(\cos\varphi, \sin\varphi)$, and let $u = 1/r$. Level~1
derives the orbit. With $\sigma = \operatorname{sign} L_z$, $L = \abs{L_z}$,
$\theta = \sigma(\varphi - \varphi_0)$ and $u'_0 = -p_{r,0}/L$,
\begin{equation}\label{eq:orbit}
  u(\theta) = B + (u_0 - B)\cos\Gamma\theta + \frac{u'_0}{\Gamma}\sin\Gamma\theta,
\end{equation}
where $\Gamma^2 = 1 - \kappa^2/(c^2 L^2)$ and $B = -\kappa W/(c^2 L^2 \Gamma^2)$, and $W$ is
the energy \xeqref{foundations}{eq:energy}. This is \xeqref{first_bend}{eq:general}.
Differentiating it gives
\begin{equation}\label{eq:uprime}
  u'(\theta) = -\Gamma(u_0 - B)\sin\Gamma\theta + u'_0\cos\Gamma\theta.
\end{equation}
The time element and the momentum along the same orbit are \xeqref{first_bend}{eq:dtdth}
and \xeqref{first_bend}{eq:momentum}. Here $\Gamma^2 > 0$: it falls short of $1$ by
$\val{gGap}$.

The launch sits up and to the left of the charge, and the momentum has a component toward
it. The radial momentum is $p_{r,0} = \val{prZero}$, so $u'_0 = \val{duZero}$ and the
launch is off any apsis. $L_z = \val{Lz} < 0$, so $\sigma = \val{sigma}$: the polar angle
about the charge decreases as the particle flies.

\subsection{The face}\label{sec:face}

The position on the orbit is
\begin{equation}\label{eq:position}
  x(\theta) = x_\mathrm{c} + \frac{\cos\varphi}{u(\theta)}, \qquad
  y(\theta) = y_\mathrm{c} + \frac{\sin\varphi}{u(\theta)}, \qquad
  \varphi = \varphi_0 + \sigma\theta.
\end{equation}
While $u > 0$,
\begin{equation}\label{eq:dx}
  \deriv{x}{\theta} = -\frac{\sigma\sin\varphi}{u} - \frac{\cos\varphi}{u^2}\,u',
\end{equation}
and the same differentiation gives $\dd y/\dd\theta$. A face $x = \mathrm{const}$ or
$y = \mathrm{const}$ is a root of one of those coordinates.

Which boundary comes first is read off the orbit. Walk $\theta$ from $0$ in steps of
$10^{-3}$, up to $\theta = 3$. On the chord from $\theta_k$ to $\theta_{k+1}$, in the
board's order (\xsecref{foundations}{sec:board}): if $r(\theta_{k+1}) \le R$, the flight
collides; if the point at $\theta_{k+1}$ lies outside the world, the flight leaves,
through the side the chord reaches first; if that point lies inside the detector, the
flight enters it, through the side the chord reaches first. The side is the one whose
value the chord attains at the smallest fraction of its length. The root is then polished
by Newton's method \xeqref{foundations}{eq:newton}, started at the middle of that chord,
with the derivative \eqref{eq:dx} or its partner for $y$. The face coordinate of the
accepted root is that face's value.

The reference walk meets the floor $y = y_D = \val{yDlo}$ while $x$ is between $\val{xD}$
and $\val{xDhi}$. The root of
\begin{equation}\label{eq:face}
  f(\theta) = y(\theta) - y_D = 0
\end{equation}
is the entry. The same orbit crosses the line $x = \val{xD}$ earlier, at an ordinate
still below $y_D$, so that line is outside the detector and is not an entry.

The closest approach on the orbit is the maximum of $u$ on the flight, a root of
$u' = 0$ \eqref{eq:uprime} when that root lies between the launch and the face. It is
far outside the sphere (§\ref{sec:hand}).

\subsection{The time and the radiated energy}\label{sec:time}

Along the orbit,
\[
  t(\theta) = \int_0^\theta \deriv{t}{\theta'}\,\dd\theta', \qquad
  \deriv{t}{\theta} = \frac{W - \kappa u}{c^2 L u^2},
\]
which is \xeqref{first_bend}{eq:dtdth}. Level~1 evaluates the integral in closed form when
$u'_0 = 0$ (\xeqref{first_bend}{eq:time}). Here $u'_0 \ne 0$, and the number compared
below is Simpson's rule \xeqref{foundations}{eq:simpson} on $\dd t/\dd\theta$, with
$\val{nSimp}$ panels, as in \xsecref{twin_beams}{sec:time}. Halving the panel count
estimates the error as $(S_n - S_{2n})/15$.

The radiated energy is the same rule applied to the power. The acceleration is
\xeqref{first_bend}{eq:acceleration} and the power is \xeqref{foundations}{eq:larmor}.
The integral uses $\val{nRad}$ panels; the uncertainty is the relative step-doubling error.

\subsection{The direction cone}\label{sec:cone}

The level file gives the cone as two angles in degrees, an axis $\alpha$ and a half-angle
$\delta$. The axis direction is the unit vector
$\uvec{e} = (\cos\alpha, \sin\alpha)$. Let $\vect{v}$ be the velocity at the first entry
into the detector box, and let $\beta$ be the angle from $\uvec{e}$ to $\vect{v}$, taken
in $[0, \pi]$. Then $\vect{v}\cdot\uvec{e} = \abs{\vect{v}}\cos\beta$ and
$\abs{\vect{v}\times\uvec{e}} = \abs{\vect{v}}\sin\beta$, so
\begin{equation}\label{eq:cone}
  \beta = \operatorname{atan2}\bigl(\abs{\vect{v}\times\uvec{e}},\, \vect{v}\cdot\uvec{e}\bigr).
\end{equation}
The entry counts when $\beta \le \delta$, with $\delta$ in radians. The outcome is then
Arrived. When the particle has entered the box and $\beta > \delta$, the box absorbs it
and the entry does not count: the outcome is Rejected (\physref{6.1}). A flight that
never enters the box is not tested.

Since $\vect{v} = \vect{p}/(\gamma m)$ with $\gamma m > 0$, the same $\beta$ is obtained
from $\vect{p}$. With $\uvec{e} = (1, 0)$ and $p_x > 0$,
$\beta = \operatorname{atan2}(\abs{p_y}, p_x)$. The test uses the state at the root of
\eqref{eq:face}, the moment of entry.

\section{The reference solution by hand}\label{sec:hand}

\begin{procedure}
  \item\label{st:launch} \emph{The launch} (\xeqref{foundations}{eq:launch}):
    $\gamma_0 = \val{gammaZero}$, $\abs{\vect{p}_0} = \val{pZero}$, and the direction is
    the file's, at $\val{launchAng}^\circ$ to $+x$.
  \item\label{st:orbit} \emph{The orbit} \eqref{eq:orbit}, about
    $(\val{xC}, \val{yC})$. The launch radius is $r_0 = \val{rZero}$ at
    $\varphi_0 = \val{phiZero}$, and
    $L = \abs{L_z} = \val{L}$ and
    \begin{align*}
      L_z &= \val{Lz}, & W &= \val{W}, & \Gamma &= \val{Gamma}, \\
      B &= \val{B}, & u_0 &= \val{uZero}, & u'_0 &= \val{duZero}.
    \end{align*}
  \item\label{st:face} \emph{The first event}, by the walk of §\ref{sec:face}. Five points
    of the orbit:
    \begin{center}\small\val{samples}\end{center}
    The line $x = \val{xD}$ is crossed at $\theta = \val{thetaLeft}$, where
    $y = \val{yAtLeft}$, below the floor $y_D = \val{yDlo}$. The walk then meets the
    floor while $x$ is inside the detector. Newton's method on \eqref{eq:face}, started
    at $\theta = \val{thGuess}$:
    \begin{center}\small\val{newton}\end{center}
    After $\val{newtonSteps}$ corrections the residual is $\val{newtonRes}$, and
    $\theta_\mathrm{e} = \val{thetaE}$. The entry is
    $(\val{xE}, \val{yE})$, on the floor. The closest approach is at
    $\theta = \val{thMin}$, where $r = \val{rmin}$.
  \item\label{st:time} \emph{The time and the momentum.} Simpson's rule with
    $\val{nSimp}$ panels gives $t_\mathrm{e} = \val{tE}$, uncertain by $\val{terr}$, and
    far below $t_{\max}$. The momentum \xeqref{first_bend}{eq:momentum} is
    $(\val{pxE}, \val{pyE})$. The kinetic energy, which this detector does not test, is
    $\val{TE}$.
  \item\label{st:cone} \emph{The cone} \eqref{eq:cone}. Both components of the momentum
    are positive, so $\beta = \operatorname{atan2}(\abs{p_y}, p_x) = \val{devE}$
    radians, which is $\val{angE}^\circ$. The half-angle is $\val{halfDeg}^\circ$, and
    the margin is $\val{marginE}^\circ$. The entry counts.
  \item\label{st:radiation} \emph{The radiated energy.} Simpson's rule on the power, with
    $\val{nRad}$ panels: $E_\mathrm{rad} = \val{Erad}$, uncertain by the relative amount
    $\val{radErr}$.
\end{procedure}

\begin{result}[The reference flight]
The particle enters the detector through the floor $y = \val{yDlo}$ at
$(\val{xE}, \val{yE})$, at $t_\mathrm{e} = \val{tE}$, with momentum
$(\val{pxE}, \val{pyE})$. The arrival direction is $\val{angE}^\circ$ from $+x$, inside
the cone of half-angle $\val{halfDeg}^\circ$. The level is solved.
\end{result}

The orbit is the closed form \eqref{eq:orbit}. The Newton table is the iteration for the
face, and the two quadratures are specified in §\ref{sec:time}; the time is uncertain by
$\val{terr}$ and the radiated energy by the relative amount $\val{radErr}$.

\section{Comparison with the game}\label{sec:comparison}

\comparisontable

The orbit is exact. The time quadratures are uncertain by $\val{terr}$ on the reference
and by $\val{halfTerr}$ on the weaker charge, both far below $10^{-12}$. The tolerance on
a state, and on the deviation $\beta$, is $10^{-9}$, ten times the accuracy of a verified
flight (\xsecref{foundations}{sec:verification}). The tolerance on a radiated energy is
$10^{-4}$ relative, and those integrals are uncertain by $\val{radErr}$ and
$\val{halfRadErr}$. Both tolerances were set from those estimates. With no charge the power is
identically zero, so that radiated energy is compared absolutely.

The reference integration stops on the floor, so the ordinates agree exactly. The straight
ray and the weaker charge stop on $x = \val{xD}$, so those abscissae agree exactly. The
largest absolute difference in the table is $\val{cmpRefTDiff}$, the reference arrival
time. It sits at the game's accuracy, above the quadrature error. The radiated energies of
the two charged flights agree to the relative differences $\val{cmpRefRadDiff}$ and
$\val{cmpHalfRadDiff}$.

\section{Other solutions}\label{sec:others}

Two further placements, each compared above. Both enter the detector. The cone is what
rejects them.

\begin{enumerate}
  \item No charge. The momentum never changes, so the ray is uniform motion at speed
    $\val{vZero}$ along the launch direction. It meets the left face $x = \val{xD}$ at
    $t = \val{freeT}$, at the point $(\val{freeX}, \val{freeY})$, before it would leave
    the world. The ordinate lies between $\val{yDlo}$ and $\val{yDhi}$, so this is an
    entry. The direction is the launch direction, $\beta = \val{freeAng}^\circ$, and the
    margin is $\val{freeMargin}^\circ$. The entry does not count.
  \item $Q = \val{halfQ}$ at $\val{halfNode}$, so $\kappa = \val{halfK}$: the reference
    point, with half the coupling. The flight is the orbit \eqref{eq:orbit}. The walk
    meets the left face while the ordinate $\val{halfY}$ is inside the detector, at
    $t = \val{halfT}$ (quadrature $\val{halfTerr}$). The momentum is $(\val{halfPx}, \val{halfPy})$, so
    $\beta = \val{halfAng}^\circ$ and the margin is $\val{halfMargin}^\circ$. The bend
    has lowered the ray from the straight flight, and the direction is still outside the
    cone. The closest approach is $\val{halfRmin}$, and the radiated energy is
    $\val{halfErad}$.
\end{enumerate}

The reference works because $\kappa = \val{kappa}$ is strong enough to drop the arrival
from $\val{freeAng}^\circ$ to $\val{angE}^\circ$, and the floor crossing at
$x = \val{xE}$ is still inside the detector. Half that coupling leaves the entry on the
left face at $\val{halfAng}^\circ$. With no charge the ray enters at the launch angle.

\section{What the game computes}\label{sec:engine}

The game integrates \xeqref{foundations}{eq:eom} with its eighth-order Runge--Kutta method
at the tolerances $10^{-10}$ and $10^{-12}$ (\physref{5}). It finds the crossing on its
interpolant (\physref{6}) and applies \eqref{eq:cone} to the velocity at that entry
(\physref{6.1}). It verifies the outcome by comparing its two runs (\physref{7}). The
model is this document's, including \ref{as:radiation} (\physref{8}). The agreement in
§\ref{sec:comparison} is between the orbit \eqref{eq:orbit}, the cone \eqref{eq:cone} and
that integration.

\section{Further reading}\label{sec:reading}

\begin{itemize}
  \item \levelref{first_bend}: the orbit \eqref{eq:orbit}, the time element and the
    momentum. There the charge is at the origin and the launch is an apsis, so the time
    is a closed form.
  \item \levelref{twin_beams}: the same orbit with $u'_0 \ne 0$, and Simpson's rule for
    the time.
  \item \xsecref{foundations}{sec:board}: a detector may require a direction. The angle
    is \eqref{eq:cone}, derived in §\ref{sec:cone}.
\end{itemize}

\levelfileappendix

\end{document}
